\section{Method}

Fig.~\ref{fig:pipeline} shows the pipeline. It has six steps, one script each.

\subsection{Overlap between corpora}

We compare every pair of the nine sources on the full cleaned corpora, before
sampling, at three levels.
\begin{itemize}
\item \textit{Exact raw}: MD5 of the text exactly as the model reads it.
\item \textit{Exact normalized}: MD5 of the body after lowercasing and
removing all whitespace. This is the check we used in our preliminary work.
\item \textit{Near duplicate}: MinHash~\cite{broder} with 128 permutations over
9-character shingles of the first 2{,}000 letters and digits of the body, with
locality sensitive hashing to find candidates. Two emails count as the same
if their estimated Jaccard similarity is at least 0.8.
\end{itemize}
We first tried 5-word shingles. They missed many copies, because the Kaggle
aggregate often deletes line breaks and glues two words together
(``cream?Isn't''). Character shingles over letters and digits ignore that. To
check the matcher, we took 400 random SpamAssassin emails that had no
normalized exact match in Kaggle and searched Kaggle by brute force. 59\% had a
Kaggle email with a true Jaccard of at least 0.8, the same rate MinHash
reported for that group, and manual reading confirmed they were the same
messages with different line breaks or a cut first word.

\subsection{Leave-one-corpus-out with decontamination}

For each of the six corpora $T$ we train on the other five and test on all of
$T$. We run this twice. In the \textit{raw} setting the training data is used
as it is. In the \textit{clean} setting we first remove every training email
that has a near duplicate in $T$. The difference in F1 between the two is the
inflation caused by leaked emails. The three test-only sets are scored with a
model trained on all six corpora, again with near duplicates of the test set
removed. As a reference we also report the usual in-corpus score, an 80/20
stratified split inside each corpus.

\subsection{Models}

\textit{Classical.} TF-IDF over word unigrams and bigrams (50{,}000 features,
minimum document frequency 2) with Logistic Regression or multinomial Naive
Bayes, as in our preliminary work.

\textit{DistilBERT.} \texttt{distilbert-base-uncased}~\cite{distilbert} fine-tuned
for one epoch, 128 tokens, batch 16, learning rate $5\times10^{-5}$, on a
stratified sample of 1{,}000 emails per training corpus (about 5{,}000 per
fold). It runs on the laptop CPU (Apple M1). Only the clean setting is run.

\textit{Small open LLMs.} Llama-3.2-1B, Llama-3.2-3B, Llama-3.1-8B,
Qwen-2.5-7B, Gemma-3-12B and Phi-4 (14B), all instruction-tuned and called
through OpenRouter. The system prompt asks whether the email is ``malicious or
unwanted (phishing, scam or spam) or legitimate'' and to answer with one word.
Temperature is 0, output is limited to 5 tokens, and emails are cut to 1{,}500
characters. An answer containing ``phish'', ``spam'' or ``malicious'' counts as
positive, one containing ``legit'' as negative. Anything else (for example a
refusal) counts as negative and is reported. An LLM never sees our training
data, so every test set is unseen for it. In the few-shot setting the prompt
holds four short examples (two positive, two negative) drawn from corpora
other than the test corpus, which keeps the leave-one-corpus-out rule.

\subsection{Metrics and cost}

We report F1 for the positive class, because accuracy hides the failure we care
about (a model that calls everything legitimate still gets high accuracy on a
mostly legitimate corpus). The headline robustness number is the mean F1 over
the six held-out corpora, with a 95\% interval from 1{,}000 bootstrap
resamples of the emails in each test set, plus the worst of the six. For
LLMs, cost is the sum of the \texttt{usage.cost} field that OpenRouter returns
for every call, divided by the number of emails. Local models have no API
cost. We report their inference time instead.
